\documentclass[10pt,a4paper]{amsart}
\usepackage[margin=19mm]{geometry}
\usepackage{amsmath,amssymb,mathtools}
\usepackage[scr=rsfs,cal=euler]{mathalfa}
\usepackage{xfrac}
\usepackage[unicode,hidelinks,bookmarks=false]{hyperref}
\newcommand{\ie}{i.e.\ }
\newcommand{\cf}{cf.\ }
\newcommand{\resp}{respectively\ }
\newcommand{\Subsection}{\subsection}
\providecommand{\nobreakdash}{\nobreak\hbox{-}}
\setlength{\emergencystretch}{2em}
\multlinegap0pt
\usepackage{bm}
\usepackage{bbm}
\usepackage{dsfont}
\usepackage{xspace}
\usepackage{cancel}
\usepackage{mathtools}
\usepackage{amsthm}
\makeatletter
\@ifundefined{thm}{\newtheorem{thm}{Thm}}{}
\@ifundefined{lem}{\newtheorem{lem}[thm]{Lemma}}{}
\makeatother
\def\supp{{\mathop\mathrm{\,supp\,}}}
\title[Improved Berezin-Li-Yau inequality and Kr\"oger inequality a]{Improved Berezin-Li-Yau inequality and Kr\"oger inequality and consequences}
\author{Zaihui Gan, Renjin Jiang, Fanghua Lin}
\date{Source: arXiv:2507.20330v2 (CC BY 4.0)}
\begin{document}
\maketitle

\noindent\emph{Source and licence.} Improved Berezin-Li-Yau inequality and Kr\"oger inequality and consequences. Version arXiv:2507.20330v2 (CC BY 4.0), distributed under the Creative Commons Attribution 4.0 International licence, as recorded in the arXiv licence metadata for this exact version and shown on the abstract page https://arxiv.org/abs/2507.20330v2. Full rights notice: licenses/arxiv-2507.20330-CC-BY-4.0.txt.\par\smallskip

\noindent\textit{Editorial scope.} Counted unit: the Lem whose source label is \texttt{lem-3} in the article (source0.tex lines 584--594, source label \texttt{lem-3}), delivered together with the article's own complete original proof of it (source0.tex lines 595--655, one \texttt{proof} environment of 61 lines). No mathematics is written, repaired or re-derived: statement and proof are the article's own text line for line. Required context retained uncounted: none. Every definition, display and label of those lines is retained unchanged. Omitted from this package and not redistributed: 1--583, 656--1562. Nothing inside a retained excerpt is omitted or altered: every retained body is delivered exactly as the article's lines are written, so no omission receipt of any kind accompanies this package. Citations: the retained excerpts carry no citation command at all, so no bibliography entry of the article is transcribed and no reference is redistributed. Reference adaptation: none was needed and none was made; every reference inside the delivered unit resolves inside it, so no unresolved reference or citation is produced. Printed slips kept literal: no printed word, symbol, hypothesis or number of the counted statement or of its proof was altered. Layout and class: the article's own class is \texttt{article} with packages including amssymb, amsmath, amsfonts, amsthm, color, mathrsfs, upgreek, txfonts, tocbibind, srcltx, citeref, hyperref, xcolor, graphicx; the wrapper is the lane's pinned \texttt{amsart} editorial wrapper at 10pt on A4, which restates the theorem-like environments the retained text needs and adds the article's own macro definitions that the retained text needs (\texttt{\textbackslash supp}), with extra wrapper packages (bm, bbm, dsfont, xspace, cancel, mathtools). The delivered numbering is the unit's own. Assets: no retained span contains a figure environment, a graphics-inclusion command, a PDF-page inclusion, a table or a TikZ picture; no image file is copied into this package, the package declares no asset dependency and no image file is redistributed with it. Licence: version arXiv:2507.20330v2 of this article is distributed under the Creative Commons Attribution 4.0 International licence, as recorded in the arXiv licence metadata for this exact version and shown on the abstract page https://arxiv.org/abs/2507.20330v2. A full rights notice is delivered at licenses/arxiv-2507.20330-CC-BY-4.0.txt. Characters: every retained excerpt is pure ASCII, so the delivered engine, XeLaTeX with its default Latin Modern text font, reports no missing character.
\begin{lem}\label{lem-3}
Let $n\ge 2$ and $\lambda>0$. Suppose that $g\ge 0$ on $(0,\infty)$, which satisfies
$$\int_0^\infty g(r)\,dr\le A,\, g(r)\le B$$
with $A,B>0$ satisfying $A/B<\sqrt{\frac{n+2}{n}\lambda}$.
Then
$$\int_0^\infty \left(\frac{\lambda r^n}{n} -\frac{r^{n+2}}{n+2}\right)g(r)\,dr\le B\left(\frac{\lambda \left[(t+A/B)^{n+1}-t^{n+1}\right]}{n(n+1)}-\frac{\left[(t+A/B)^{n+3}-t^{n+3}\right]}{(n+2)(n+3)}\right),$$
where  $t<\sqrt\lambda$  is such that $t+A/B>\sqrt{\lambda}$ and
\begin{align*}
\frac{\lambda \left[(t+A/B)^{n}-t^{n}\right]}{n}=\frac{\left[(t+A/B)^{n+2}-t^{n+2}\right]}{(n+2)}.
\end{align*}
\end{lem}

\begin{proof}
Note that the function
$$\frac{\lambda r^n}{n} -\frac{r^{n+2}}{n+2}$$
tends to zero when $r\to 0$, and is minus infinite when $r=\infty$.
Moreover, there is only one maximum in $(0,\infty)$, which attains at the point
$$r=\sqrt\lambda,$$
and the function is negative if
$$\lambda< \frac{n}{n+2}r^2\Longleftrightarrow r>\sqrt{\frac{n+2}{n}\lambda}.$$
So for any $g$ satisfies the assumption,
$$\int_0^\infty \left(\frac{\lambda r^n}{n} -\frac{r^{n+2}}{n+2}\right)g(r)\,dr\le \int_0^{\sqrt{\frac{n+2}{n}\lambda}} \left(\frac{\lambda r^n}{n} -\frac{r^{n+2}}{n+2}\right)g(r)\,dr.$$
Therefore we may and do assume that $\supp g\subset (0,\sqrt{\frac{n+2}{n}\lambda})$.

Further, note that $\sqrt{\frac{n+2}{n}\lambda}-A/B>0$.
If $\int_0^{\sqrt{\frac{n+2}{n}\lambda}}g\,dr<A$, we can find another $\tilde g$ such that
$B\ge \tilde g\ge g$ and $\int_0^{\sqrt{\frac{n+2}{n}\lambda}}\tilde g\,dr=A$. It is obvious that
$$\int_0^{\sqrt{\frac{n+2}{n}\lambda}} \left(\frac{\lambda r^n}{n} -\frac{r^{n+2}}{n+2}\right)g(r)\,dr\le \int_0^{\sqrt{\frac{n+2}{n}\lambda}} \left(\frac{\lambda r^n}{n} -\frac{r^{n+2}}{n+2}\right)\tilde g(r)\,dr.$$


For $t\in (0,\sqrt{\frac{n+2}{n}\lambda}-A/B )$, let us set
$$g_t(r)=
\begin{cases}
B, & r\in (t,t+A/B)\\
0, & \mbox{otherwise}.
\end{cases}
$$
Then
\begin{align*}
\int_0^{\sqrt{\frac{n+2}{n}\lambda}} \left(\frac{\lambda r^n}{n} -\frac{r^{n+2}}{n+2}\right)g_t(r)\,dr&=B\left(\frac{\lambda r^{n+1}}{n(n+1)}-\frac{r^{n+3}}{(n+2)(n+3)}\right)\bigg|_{t}^{t+A/B}\\
&=B\left(\frac{\lambda \left[(t+A/B)^{n+1}-t^{n+1}\right]}{n(n+1)}-\frac{\left[(t+A/B)^{n+3}-t^{n+3}\right]}{(n+2)(n+3)}\right)\\
&=:B G(t).
\end{align*}
The last term $G(t)$ as a function of $t$, converges to $-\infty$ as $t\to\infty$, and when $t=0$, it holds that
\begin{align*}
\left(\frac{\lambda \left[(A/B)^{n+1}\right]}{n(n+1)}-\frac{\left[(A/B)^{n+3}\right]}{(n+2)(n+3)}\right)& =(A/B)^{n+1} \left(\frac{\lambda }{n(n+1)}-\frac{\left[(A/B)^{2}\right]}{(n+2)(n+3)}\right)\\
&\ge (A/B)^{n+1} \left(\frac{\lambda }{n(n+1)}-\frac{\lambda}{n(n+3)}\right)\\
&= (A/B)^{n+1} \frac{2\lambda }{n(n+1)(n+3)}.
\end{align*}
The maximal point of the $G(t)$ is attained at the point $t_0$ such that
\begin{align*}
\frac{\lambda \left[(t_0+A/B)^{n}-t_0^{n}\right]}{n}=\frac{\left[(t_0+A/B)^{n+2}-t_0^{n+2}\right]}{(n+2)}.
\end{align*}
Equivalently,
\begin{align*}
\frac{\lambda(t_0+A/B)^{n}}{n}-\frac{\lambda(t_0+A/B)^{n+2}}{n+2} =\frac{\lambda t_0^{n}}{n}-\frac{t_0^{n+2}}{(n+2)}.
\end{align*}
Since $\frac{\lambda r^n}{n} -\frac{r^{n+2}}{n+2}$ is increasing on $(0,\sqrt\lambda)$ and decreasing on $(\sqrt\lambda,\infty)$, we see that
$t_0<\sqrt\lambda$ and $t_0+A/B>\sqrt{\lambda}$.


For any other $\supp g\subset (0,\sqrt{\frac{n+2}{n}\lambda})$ satisfying $\int_0^{\sqrt{\frac{n+2}{n}\lambda}} g\,dr=A$ and the assumptions,
suppose that $g$ is not fully supported in $[t_0,t_0+A/B]$, i.e.,
$$\int_{t_0}^{t_0+A/B} g(t)\,dt<A.$$
Then
\begin{align*}
\int_{(0,\sqrt{\frac{n+2}{n}\lambda})\setminus [t_0,t_0+A/B]} \left(\frac{\lambda r^n}{n} -\frac{r^{n+2}}{n+2}\right)g(r)\,dr&\le
\sup_{t\in (0,\sqrt{\frac{n+2}{n}\lambda})\setminus [t_0,t_0+A/B]}\left(\frac{\lambda r^n}{n} -\frac{r^{n+2}}{n+2}\right) \left(A-\int_{t_0}^{t_0+A/B} g(t)\,dt\right) \\
&=\inf_{t\in [t_0,t_0+A/B]}\left(\frac{\lambda r^n}{n} -\frac{r^{n+2}}{n+2}\right) \left(\int_{t_0}^{t_0+A/B} \left(g_{t_0}(t)-g(t)\right)\,dt\right)\\
&\le \int_{t_0}^{t_0+A/B} \left(\frac{\lambda r^n}{n} -\frac{r^{n+2}}{n+2}\right)\left(g_{t_0}(t)-g(t)\right)\,dt.
\end{align*}
The proof is complete.
\end{proof}

\end{document}
